% TOP_PROBLEM: {"id": 3225, "title": "Strong colorability", "category_id": 3, "category": {"id": 3, "name": "graph_theory", "display_name": "Graph Theory", "description": "Problems involving graphs, networks, and their properties.", "slug": "graph-theory", "order_index": 3, "created_at": "2026-07-31T15:26:25.671Z"}, "status": "open", "problem_number": "OPG-171", "set_id": 12, "statusQualification": "The substantive 2 Delta assertion is stated for Delta at least one; edgeless graphs are handled separately to avoid a zero-colour convention."}
% FIRST_POSTED: 2026-10-01
% ENTRY_KIND: statement_only
% UnsolvedMath ID 3225; source catalogue code OPG-171
% !TeX program = lualatex
% Definition and sources only; this file is not an LLM solution attempt.
\documentclass[11pt,a4paper]{article}
\usepackage[margin=25mm]{geometry}
\usepackage{fontspec}
\setmainfont{lmroman10-regular.otf}[BoldFont=lmroman10-bold.otf,
 ItalicFont=lmroman10-italic.otf,BoldItalicFont=lmroman10-bolditalic.otf]
\usepackage{amsmath,amssymb,mathtools}
\usepackage{unicode-math}
\setmathfont{latinmodern-math.otf}
\AtBeginDocument{\let\setminus\smallsetminus}
\usepackage{xurl,hyperref}
\hypersetup{colorlinks=true,urlcolor=blue}
\setcounter{secnumdepth}{0}
\setlength{\parindent}{0pt}
\setlength{\parskip}{5pt}
\setlength{\emergencystretch}{3em}
\begin{document}
\section*{Strong colorability}\label{3225}
{\small UnsolvedMath ID 3225 · OPG-171 · Graph Theory · OpenGarden}
\subsection{Definitions and mathematical statement}
Let $G=(V,E)$ be a finite
simple graph and let $r\geq1$
be an integer. A partition
$\mathcal P$ of $V$ into
parts of size at most $r$
consists of nonempty pairwise
disjoint subsets whose union
is $V$. A strong $r$-colouring
relative to $\mathcal P$ is
a function $c:V\to\{1,\ldots,r\}$
that is proper on $G$ and
injective on every part of
$\mathcal P$.

The graph is strongly
$r$-colourable if such a
colouring exists for every
partition into parts of size
at most $r$. Different
partitions may require different
colourings. Equivalently, add
all edges within each part
to make it a clique; every
graph produced this way must
have chromatic number at
most $r$.

The conjecture states that
every finite simple graph
of maximum degree $\Delta\geq1$
is strongly $2\Delta$-colourable.
Explicitly, after any partition
into parts of size at most
$2\Delta$, one must colour
the vertices with $2\Delta$
colours while distinguishing
both adjacent vertices and
vertices sharing a part.

The maximum degree is measured
in $G$ before adding the
partition edges. There is
no restriction on how parts
interact with edges, and no
assumption that $|V|$ is
divisible by $2\Delta$.
The positive-degree condition
avoids the undefined phrase
strongly zero-colourable; edgeless
graphs are trivially strongly
$r$-colourable for every
positive $r$ under this
definition.
\subsection{Short English statement}
After any partition into groups of size at most twice the maximum degree, can a graph be coloured with that many colours so every group also receives distinct colours?
\subsection{Sources}
\begin{itemize}
\item[S1] ULAM.AI and the UnsolvedMath Contributors, supplied problems.json, record 3225 (OPG-171), OpenGarden collection. Catalogue identity and source transcription; CC BY 4.0.
\url{https://huggingface.co/datasets/ulamai/UnsolvedMath}.
\item[S2] Open Problem Garden, Strong colorability. Original problem and discussion; the mathematical formulation below is expanded from the supplied source record.
\url{http://www.openproblemgarden.org/op/strong_colorability}.
\end{itemize}
\end{document}
